\begin{tikzpicture}[scale=0.9,
  abierto/.style={draw=black, thick, dashed, rounded corners=5pt}
]
  % Contraejemplo: topología generada por los n_>= = {m in N | m >= n}
  % Columna k: U_k y, dentro, U_{k+1}
  \definecolor{Crema}{HTML}{FEFAE0}
  \definecolor{Oliva}{HTML}{C5C9B6}
  \definecolor{Rosa}{HTML}{E9C1CD}
  \colorlet{Azul}{FancyBlue!25}
  \useasboundingbox (-0.6,-4.0) rectangle (4.3,0.95);
  \def\paso{0.6}
  \def\fondo{-3.35}
  \foreach \k/\cext/\cint/\s in {0/white/Crema/8, 1/Crema/Oliva/9, 2/Oliva/Rosa/10, 3/Rosa/Azul/11}{%
    \pgfmathsetmacro{\x}{\k*1.15}
    \pgfmathsetmacro{\te}{-\k*\paso+0.3}
    \pgfmathsetmacro{\ti}{-(\k+1)*\paso+0.27}
    \pgfmathtruncatemacro{\kk}{\k+1}
    \uncover<\s->{%
      % U_k (exterior)
      \fill[\cext, rounded corners=5pt] (\x-0.36,\fondo) -- (\x-0.36,\te) -- (\x+0.36,\te) -- (\x+0.36,\fondo) -- cycle;
      % U_{k+1} (interior)
      \fill[\cint, rounded corners=4pt] (\x-0.22,\fondo) -- (\x-0.22,\ti) -- (\x+0.22,\ti) -- (\x+0.22,\fondo) -- cycle;
      \draw[abierto] (\x-0.36,\fondo) -- (\x-0.36,\te) -- (\x+0.36,\te) -- (\x+0.36,\fondo);
      \draw[abierto, rounded corners=4pt] (\x-0.22,\fondo) -- (\x-0.22,\ti) -- (\x+0.22,\ti) -- (\x+0.22,\fondo);
      \foreach \n in {\k,...,5}
        \node at (\x,-\n*\paso) {$\n$};
      \node at (\x,\te+0.3) {$U_{\k}$};
      \draw[->] (\x,\fondo-0.08) -- (\x,\fondo-0.35);
    }
  }
  % Continúa a la derecha
  \uncover<11->{%
    \foreach \n in {3,4,5}
      \draw[->] (3.9,-\n*\paso) -- (4.2,-\n*\paso);
  }
\end{tikzpicture}
